\section{Semantic Session Typing for Deadlock Freedom}
\label{sec:logicalRelation}

The semantic typing judgment $\senv \proves_{\wrld} \proc : \ty$ of this
section fulfils three goals.
\begin{enumerate}[(1)]
  \item It ensures deadlock freedom for closed processes (``adequacy'');
  \item
    It validates all rules of the static system of
    \cref{sec:syntacticTyping}, which gives the ``fundamental property''
    that
    $
      \senv \typesW{\wrld} \proc : \ty
      \implies
      \senv \proves_{\wrld} \proc : \ty
    $;
  \item It is defined on untyped processes and constrains a process only
    at its external interface, so a safe process that the static system
    refuses can belong to it.
\end{enumerate}
Goal (3) is a property of the definition. No example in this paper
establishes membership for a process that the static system refuses, and
every example of \cref{sec:examples} goes through
\cref{th:fundamental}.


\subsection{Worlds and Readiness}
\label{sec:worlds}

The relation is indexed by a \emph{world}, which carries the order the
points of a running program stand in.

\begin{definition}[World \rocq{world}, \rocq{world_wf}, \rocq{strict}]
\label{def:world}
  A world $\wrld = (N,{\prec}_{\wrld})$ consists of a finite set of
  names $N$ and a relation $\fLt{}{}$ on points.
  Its \emph{support} is $\supp(\wrld) := N$. It is
  \emph{well formed} when $\fLt{}{}$ is transitive and contains program
  order $\fLt{\pntA}{\poff{\pntA}{1}}$ and
  $\fLt{\pntA}{\poff{\pdual{\pntA}}{1}}$, every support name is a point
  constant, and outside the support only program order holds. It is \emph{strict} when it is well formed and $\fLt{}{}$ is
  irreflexive.
  $\wrld$ \emph{knows} $\pntA$ when $\rootof(\pntA)$ is in its support
  (\rocq{pt_known}).
\end{definition}

\begin{definition}[Extension \rocq{wext}, \rocq{wext_fresh}]
\label{def:wext}
  $\wrld \wextsym \wrld'$ holds when $\wrld'$ is strict and contains the
  support and the order of $\wrld$.
  $\wrld \wextfresh \wrld'$ \emph{placing $N$} holds in addition when
  $N$ is disjoint from the support of $\wrld$, the support of $\wrld'$
  is that of $\wrld$ with $N$, and $\wrld'$ adds no fact between points
  whose roots avoid $N$.
\end{definition}

Only $\tNew$ extends the world. A new channel either sits at a point the
world knows and leaves the world unchanged, or it extends the world
freshly (\rocq{new_place}). A creation
constraint is exactly what a fresh extension can realize: if $\cstr$ is a
creation constraint under the facts of a strict $\wrld$ over points
$\wrld$ knows, then some $\wrld \wextfresh \wrld'$ placing a fresh
constant $\pcst$ satisfies $\tysub{\cstr}{\pcst}{\pviA}$
(\rocq{wext_cc}). The two
results below are what makes a composed run strict. Write
$\supp(\wrld)$ for the support of $\wrld$, write
$\wrld_1 \wjoinsym \wrld_2$ for the world on
$\supp(\wrld_1) \cup \supp(\wrld_2)$ generated by the orders of both
(\rocq{wjoin}), and write $\wrest{\wrld}{A}$ for the world on
$\supp(\wrld) \cap A$ that keeps the facts of $\wrld$ whose two points
have their roots in $A$ (\rocq{wrestrict}).

\begin{theorem}[Amalgamation \rocq{amalgamation}]
\label{thm:amalgamation}
  For all strict $\wrld_1$ and $\wrld_2$, and for
  $A = \supp(\wrld_1) \cap \supp(\wrld_2)$,
  \[
    \strict(\wrld_1 \wjoinsym \wrld_2)
    \iff
    \strict(\wrest{\wrld_1}{A} \wjoinsym \wrest{\wrld_2}{A}) .
  \]
\end{theorem}

(\rocq{wjoin_wf}, \rocq{wrestrict_wf}, \rocq{wjoin_restrict_incl}, \rocq{jpath},
\rocq{join_jpath}, \rocq{jpath_decomp}, \rocq{join_cycle}).

\begin{theorem}[Fresh extension \rocq{fresh_extension}]
\label{thm:fresh-extension}
  If $\wrld_1 \wextfresh \wrld_1'$ placing $N$, the support of
  $\wrld_2$ avoids $N$, $\wrld_2$ is strict and
  $\wrld_1 \sqcup \wrld_2$ is strict, then
  $\wrld_1' \sqcup \wrld_2$ is strict.
\end{theorem}

(\rocq{wext_fresh}, \rocq{wgen_incl}).

The step clauses of the relation are gated by \emph{readiness}. A term
or process with no internal step left must either be done or offer some
of the entries it holds, and the offered entries must lie strictly below
the partner points of the entries it holds but does not offer. Entries
are compared, not names. A thread that holds both ends of $\ch$, at
$\pntA$ and $\pdual{\pntA}$, and blocks on the first holds the second
unoffered and would need $\fLt{\pntA}{\pntA}$.

\begin{definition}[Readiness \rocq{ready_t}, \rocq{ready_p}]
\label{def:ready}
  $\ents(\tenv)$ is the set of entries of $\tenv$, one per end,
  $\hd(\entry)$ is the point of the entry $\entry$ (\rocq{entry_head}),
  and $O$ ranges over sets of entries. The two conditions on $O$ are
  \begin{align*}
    \belowrestop_{\wrld}(\tenv, O, \frm)
      &\iff
      \forall \entry \in O \cap \ents(\tenv) \st
        \forall \entry' \in \ents(\tenv) \setminus O \st
          \wrld \types \fLt{\hd(\entry)}{\pdual{\hd(\entry')}}
    \\
      &\qquad \text{and }
      \forall \entry \in O \cap \ents(\tenv) \st
        \forall \pntB \in \frm \st
          \wrld \types \fLt{\hd(\entry)}{\pntB}
    \\
    \aboveofferedop_{\wrld}(\tenv, O)
      &\iff
      \forall \entry \in \ents(\tenv) \setminus O \st
        \exists \entry' \in O \cap \ents(\tenv) \st
          \wrld \types \fLt{\hd(\entry')}{\pdual{\hd(\entry)}}
  \end{align*}
  (\rocq{below_rest}, \rocq{above_offered}). Let $\offered(\tm)$ and
  $\offered(\proc)$ be the entries on which $\tm$ and $\proc$ offer a
  session label (\rocq{offered_t}, \rocq{offered_p}).
  \begin{itemize}
    \item $\readyop_{\wrld}(\tenv \mid \frm, \tm)$ holds when $\tm$
      having no $\beta$, $\lNew{\ch}$ or $\lSpawnL{\val}$ step implies
      \[
        \big(\offered(\tm) = \emptyset \implies \tm \text{ is a value}\big)
        \land
        \big(\tm \text{ has a forwarding redex} \lor
          \belowrestop_{\wrld}(\tenv, \offered(\tm), \frm)\big)
      \]
    \item $\readyop_{\wrld}(\senv, \proc)$ holds when $\proc$ having no
      $\tau$ step implies that $\proc$ has a forwarding redex, or
      \[
        \big(\offered(\proc) = \emptyset \implies
          \proc \text{ is terminated}\big)
        \land
        \aboveofferedop_{\wrld}(\tenvof{\senv}, \offered(\proc))
      \]
  \end{itemize}
\end{definition}

The two clauses are asymmetric. A term owes $\belowrestop$, which puts
every offered entry below the partner point of every unoffered one and
below every point of the future. A process owes only $\aboveofferedop$,
which puts the partner point of every unoffered entry above some offered
one, and a process has no future. A
forwarding never blocks, so a term or process at a forwarding redex
compares nothing, and only the process clause lets a forwarding redex
discharge the condition on values as well.

\begin{example}[Ready and blocked threads]
  Let $\wrld$ be strict with $\supp(\wrld)=\set{\pcst,\asTy{b}}$ and
  $\wrld\types\fLt{\pcst}{\pdual{\asTy{b}}}$, and
  $\wrld\not\types\fLt{\pdual{\asTy{b}}}{\pcst}$. In
  \[
    (\set{\ch,\chB},\mset{
      \tSeq{\tApp{\tClos}{\ch}}{\tApp{\tClos}{\chB}},
      \tSeq{\tApp{\tClos}{\chB}}{\tApp{\tClos}{\ch}}})
  \]
  the first thread holds $\ch:\sAt{\sTerm}{\pcst}$ and
  $\chB:\sAt{\sTerm}{\asTy{b}}$, the second $\ch:\sAt{\sTerm}{\pdual{\pcst}}$
  and $\chB:\sAt{\sTerm}{\pdual{\asTy{b}}}$.
  At empty future the first thread is ready because its offer at
  $\pcst$ precedes the partner point $\pdual{\asTy{b}}$ of its unoffered
  end. The second thread is not ready because its offer at
  $\pdual{\asTy{b}}$ would need $\fLt{\pdual{\asTy{b}}}{\pcst}$.
  Neither close can synchronize, so the process is blocked.
\end{example}

\begin{example}[A future delays readiness]
  In the same world, $\tApp{\tClos}{\ch}$ with
  $\tenv=\ch:\sAt{\sTerm}{\pcst}$ is ready at future
  $\set{\pdual{\asTy{b}}}$. It is not ready at future $\set{\pcst}$ because
  $\belowrestop$ would require $\wrld\types\fLt{\pcst}{\pcst}$.
\end{example}

\subsection{A Logical Relation for \OurLang}
\label{sec:logicalRelation:relation}

Three relations express the typing constraints
(\cref{fig:logrel-rules-value,fig:logrel-rules-term,fig:logrel-rules-proc}):
$\lrVW{\wrld}{n}{\tenv \types \ty}$ for values,
$\lrEW{\wrld}{n}{\tenv \types \ty}{\frm}$ for terms, and
$\lrPW{\wrld}{n}{\senv \types \ty}$ for processes.
The index $n \in \Nat$ is a step index~\cite{Ahmed06} bounding the
number of transitions the relation examines.
$\senv$ and $\tenv$ are two types, not two names for one thing. A
session environment $\senv$ maps each channel to one end
$\sAt{\sty}{\pntA}$, a term environment $\tenv$ maps each channel to one
end or to both, and $\tenvof{\senv}$ turns every entry of $\senv$ into a
single end (\rocq{senv}, \rocq{tenv}, \rocq{ends}, \rocq{tenv_of}). A
component's interface exposes one end of each channel it talks on, so a
process reads $\senv$. The relation must also describe a channel whose
two ends are owned inside the component, after a cut or straight after
$\tNew$, where \ref{rule:O-New} installs the two ends at dual points, so
a term reads $\tenv$. $\heads(\tenv)$ is the set of partner points of
the entries of $\tenv$ (\rocq{tenv_heads}). $\tenv_1 \cdot \tenv_2$
splits $\tenv$, and two single ends of one
channel join when they are dual at dual points (\rocq{tenv_comp}).

\begin{figure}[tp]
  \adjustfigure[]
  \rulesection*{Value logical relation $\val \in \lrVW{\wrld}{n}{\tenv \types \ty}$ (\rocq{Vdef})}
  \begin{proofrules}
    \infer*[lab=V-Unit]{}{
      \tUnit \in \lrVW{\wrld}{n}{\tEmpty \types \tyUnit}
    } \label{rule:V-Unit}
    \and
    \infer*[lab=V-Chan]{
      \sty_0 \approx \sty
    }{
      \ch \in \lrVW{\wrld}{n}{\ch : \sAt{\sty_0}{\pntA} \types \sAt{\sty}{\pntA}}
    } \label{rule:V-Chan}
    \and
    \infer*[lab=V-Const]{
      \cst : \ty_0 \text{ an instance of a signature under } \wrld \\
      \ty_0 \approx \ty
    }{
      \cst \in \lrVW{\wrld}{n}{\tEmpty \types \ty}
    } \label{rule:V-Const}
    \and
    \infer*[lab=V-Pair]{
      \val_1 \in \lrVW{\wrld}{n}{\tenv_1 \types \ty} \\
      \val_2 \in \lrVW{\wrld}{n}{\tenv_2 \types \tyB}
    }{
      \tPair{\val_1}{\val_2} \in \lrVW{\wrld}{n}{\tenv_1 \cdot \tenv_2 \types \lPair{\ty}{\tyB}}
    } \label{rule:V-Pair}
    \and
    \infer*[lab=V-Inj]{
      \val \in \lrVW{\wrld}{n}{\tenv \types \ty_{\asTy{i}}}
    }{
      \tInj{\asTm{i}}{\val} \in \lrVW{\wrld}{n}{\tenv \types \lSum{\ty_1}{\ty_2}}
    } \label{rule:V-Inj}
    \and
    \infer*[lab=V-Abs]{
      \fchans(\tm) = \dom(\tenv) \\
      \wrld \types \fLe{\lbd}{\heads(\tenv)}
      \\\\
      \forall \wrld' \mathrel{\sqsupseteq} \wrld,\ j < n,\ \tenv'',\
        \val \in \lrVW{\wrld'}{j}{\tenv' \types \ty},\
        \frm \text{ with } \wrld' \types \fLt{\ubd}{\frm}
      \\
      \text{and } \tenv'' = \tenv \cdot \tenv' \st
        \tSub{\tm}{\val}{\vr} \in \lrEW{\wrld'}{j}{\tenv'' \types \tyB}{\frm}
    }{
      \tAbs{\vr}{\tm} \in \lrVW{\wrld}{n}{\tenv \types \ty \lolli{\ubd}{\lbd} \tyB}
    } \label{rule:V-Abs}
    \and
    \infer*[lab=V-Box]{
      \fchans(\tm) = \emptyset
      \\\\
      \forall \wrld' \mathrel{\sqsupseteq} \wrld,\ \pntvec \text{ known to } \wrld'
        \text{ with } \wrld' \types \tysub{\cstr}{\pntvec}{\pvivec} \st
      \\
      \tAbs{\tUnit}{\tSub{\tm}{\tBoxP{\vr}{\tm}}{\vr}} \in
        \lrVW{\wrld'}{n}{\tEmpty \types \tysub{(\tyUnit \lolli{\ubd}{} \ty)}{\pntvec}{\pvivec}}
    }{
      \tBoxP{\vr}{\tm} \in \lrVW{\wrld}{n}{\tEmpty \types \tyBox{\pvivec}{\cstr}{\tyUnit \lolli{\ubd}{} \ty}}
    } \label{rule:V-Box}
    \and
    \infer*[lab=V-Ex]{
      \wrld \types \fLe{\lbd}{\heads(\tenv)} \\
      \pntvec \text{ known to } \wrld \\
      \wrld \types \tysub{\cstr}{\pntvec}{\pvivec} \\
      \val \in \lrVW{\wrld}{n}{\tenv \types \tysub{\ty}{\pntvec}{\pvivec}}
    }{
      \val \in \lrVW{\wrld}{n}{\tenv \types \tyEx{\lbd}{\pvivec}{\cstr}{\ty}}
    } \label{rule:V-Ex}
  \end{proofrules}
  \rulesectionend
  \caption{The value logical relation}
  \Description{The seven rules of the value relation.}
  \label{fig:logrel-rules-value}
\end{figure}

\begin{figure}[tp]
  \adjustfigure[]
  \rulesection*{Term logical relation $\tm \in \lrEW{\wrld}{n}{\tenv \types \ty}{\frm}$ (\rocq{Edef})}
  \begin{proofrules}
    \infer*[lab=E-Zero]{}{
      \tm \in \lrEW{\wrld}{0}{\tenv \types \ty}{\frm}
    } \label{rule:E-Zero}
    \and
    \infer*[lab=E-Intro]{
      \irred \tm \implies \tm \in \lrVW{\wrld}{n+1}{\tenv \types \ty}
      \\
      \neg\irred \tm \implies \readyop_{\wrld}(\tenv \mid \frm, \tm)
      \\
      \forall \lambda, \tm'\st \tm \trans{\lambda} \tm' \implies
        \obl{\wrld ; \tenv \mid \frm}{n+1}{\tm \trans{\lambda} \tm' : \ty}
    }{
      \tm \in \lrEW{\wrld}{n+1}{\tenv \types \ty}{\frm}
    } \label{rule:E-Intro}
  \end{proofrules}
  \rulesection{Obligations of a term transition $\obl{\wrld ; \tenv \mid \frm}{n}{\tm \trans{\lambda} \tm' : \ty}$ (\rocq{Eobl})}
  \begin{proofrules}
    \infer*[lab=O-Beta]{
      \tm' \in \lrEW{\wrld}{n \cost{\lambda}}{\tenv \types \ty}{\frm}
    }{
      \obl{\wrld ; \tenv \mid \frm}{n}{\tm \trans{\beta} \tm' : \ty}
    } \label{rule:O-Beta}
    \and
    \infer*[lab=O-New]{
      \ch \notin \dom(\tenv) \implies
      \exists\, \pntA, \sty, \wrld' \st
        \wrld' \text{ places } \pntA \text{ for } \wrld, \tenv, \ty, \frm
      \\
      \tm' \in \lrEW{\wrld'}{n \cost{\lambda}}{\tenv , \ch : \sAt{\sty}{\pntA} \cdot \sAt{\dual{\sty}}{\pdual{\pntA}} \types \ty}{\frm}
    }{
      \obl{\wrld ; \tenv \mid \frm}{n}{\tm \trans{\lNew{\ch}} \tm' : \ty}
    } \label{rule:O-New}
    \and
    \infer*[lab=O-Spawn]{
      \val \in \lrVW{\wrld}{n}{\tenv_1 \types \tyUnit \lolli{\ubd}{\lbd} \tyUnit} \\
      \tm' \in \lrEW{\wrld}{n \cost{\lambda}}{\tenv_2 \types \ty}{\frm}
    }{
      \obl{\wrld ; \tenv_1 \cdot \tenv_2 \mid \frm}{n}{\tm \trans{\lSpawnL{\val}} \tm' : \ty}
    } \label{rule:O-Spawn}
    \and
    \infer*[lab=O-Recv]{
      \forall \wrld' \mathrel{\sqsupseteq} \wrld,\ j \le n \cost{\lambda},\ \tenv_2,\
        \val \in \lrVW{\wrld'}{j}{\tenv_2 \types \payloadAt{\payty}{\pntA}} \st
      \\
      \tm' \in \lrEW{\wrld'}{j}{(\tenv_1 , \ch : \sAt{\sty}{\poff{\pntA}{1}}) \cdot \tenv_2 \types \ty}{\frm}
    }{
      \obl{\wrld ; \tenv_1 , \ch : \sAt{(\sRecv{\payty}{\sty})}{\pntA} \mid \frm}{n}{\tm \trans{\lRecv{\ch}{\val}} \tm' : \ty}
    } \label{rule:O-Recv}
    \and
    \infer*[lab=O-Send]{
      \val \in \lrVW{\wrld}{n}{\tenv_2 \types \payloadAt{\payty}{\pntA}} \\
      \tm' \in \lrEW{\wrld}{n \cost{\lambda}}{\tenv_1 , \ch : \sAt{\sty}{\poff{\pntA}{1}} \types \ty}{\frm}
    }{
      \obl{\wrld ; (\tenv_1 , \ch : \sAt{(\sSend{\payty}{\sty})}{\pntA}) \cdot \tenv_2 \mid \frm}{n}{\tm \trans{\lSend{\ch}{\val}} \tm' : \ty}
    } \label{rule:O-Send}
    \and
    \infer*[lab=O-Close]{
      \tm' \in \lrEW{\wrld}{n \cost{\lambda}}{\tenv_1 \types \ty}{\frm}
    }{
      \obl{\wrld ; \tenv_1 , \ch : \sAt{\sTerm}{\pntA} \mid \frm}{n}{\tm \trans{\lTerm{\ch}} \tm' : \ty}
    } \label{rule:O-Close}
    \and
    \infer*[lab=O-Link]{
      \sty_{\chB} \isdual \sty_{\ch} \\
      \tm' \in \lrEW{\wrld}{n \cost{\lambda}}{\tenv_1 \types \ty}{\frm}
    }{
      \obl{\wrld ; \tenv_1 , \ch : \sAt{\sty_{\ch}}{\pntA} , \chB : \sAt{\sty_{\chB}}{\pdual{\pntA}} \mid \frm}{n}{}
      \\ \qquad {\tm \trans{\lLink{\ch}{\chB}} \tm' : \ty}
    } \label{rule:O-Link}
  \end{proofrules}
  \rulesectionend
  \caption{The term logical relation}
  \Description{The two rules of the term relation and the seven
    obligation rules of a term transition.}
  \label{fig:logrel-rules-term}
\end{figure}

In \ref{rule:O-New}, $\wrld'$ places $\pntA$ for $\wrld, \tenv, \ty,
\frm$ when $\pntA$ is a constant $\pcst$ fresh for all four and
$\wrld \wextfresh \wrld'$ placing $\pcst$, or when $\wrld$ knows
$\pntA$ and $\wrld' = \wrld$ (\rocq{new_place}). The first case serves
the existential typing of $\tNew$ and the second serves
$\tNewAt{\pntA}$.

A send, a close and a forwarding give a resource away, so they keep the
index. $n \cost{\lambda} = n$ for those three labels and
$n \cost{\lambda} = n - 1$ for every other (\rocq{tcost}), with
truncated subtraction. For process labels,
$n \cost{\pi} = \max(0, n - \mathit{pcost}(\pi))$ (\rocq{pcost}), except that the
$\tau$ of \ref{rule:P-Act-Link} also costs zero, so
$n \cost{\tau_{\mathrm{link}}} = n$.
The corresponding successor term is strictly smaller, which is what
makes \cref{fig:logrel-rules-term} a definition (\rocq{Efuel}).
Zero-cost process steps shrink the message size (\rocq{pcost_giving}),
so recursion at the same index is well founded.
\Cref{rule:V-Abs} reads the argument and the result at strictly smaller
indices only, so \ref{rule:V-Box} may use it at an instance of its own
type.

\begin{remark}
\label{rm:vabs}
  The side conditions on free channels of \ref{rule:V-Abs} and
  \ref{rule:V-Box} are needed because the LTS decides which channels a
  send extrudes by looking at the channels the sent value mentions
  syntactically (\cref{rm:sentc}). They and the condition on
  $\fchans(\val)$ in \ref{rule:Q-Send} are the only syntactic
  conditions in the relation. The two side conditions rule out a sent value that
  mentions channels in dead code which the continuation of the send
  still uses.
\end{remark}

\begin{figure}[tp]
  \adjustfigure[]
  \rulesection*{Process logical relation $(\chset,\pool) \in \lrPW{\wrld}{n}{\senv \types \ty}$ (\rocq{Pdef})}
  \begin{proofrules}
    \infer*[lab=P-Zero]{}{
      (\chset,\pool) \in \lrPW{\wrld}{0}{\senv \types \ty}
    } \label{rule:P-Zero}
    \and
    \infer*[lab=P-Intro]{
      \irred (\chset,\pool) \implies (\chset,\pool) = (\asPrOp{\emptyset}, \mset{\val,\tUnit,\dots,\tUnit})
        \text{ with } \val \in \lrVW{\wrld}{n+1}{\senv \types \ty}
      \\
      \neg\irred (\chset,\pool) \implies \readyop_{\wrld}(\senv, (\chset,\pool))
      \\
      \forall \pi,\chset',\pool'\st (\chset,\pool) \trans{\pi} (\chset',\pool') \implies
        \obl{\wrld ; \senv}{n+1}{(\chset,\pool) \trans{\pi} (\chset',\pool') : \ty}
    }{
      (\chset,\pool) \in \lrPW{\wrld}{n+1}{\senv \types \ty}
    } \label{rule:P-Intro}
  \end{proofrules}
  \rulesection{Obligations of a process transition $\obl{\wrld ; \senv}{n}{(\chset,\pool) \trans{\pi} (\chset',\pool') : \ty}$ (\rocq{Pobl})}
  \begin{proofrules}
    \infer*[lab=Q-Tau]{
      \chset' \disjoint \dom(\senv) \implies
      \exists\, N \text{ fresh for } \senv, \ty,\
        \wrld \wextfresh \wrld' \text{ placing } N \st
      (\chset',\pool') \in \lrPW{\wrld'}{n \cost{\pi}}{\senv \types \ty}
    }{
      \obl{\wrld ; \senv}{n}{(\chset,\pool) \trans{\tau} (\chset',\pool') : \ty}
    } \label{rule:Q-Tau}
    \and
    \infer*[lab=Q-Recv]{
      \forall \wrld' \mathrel{\sqsupseteq} \wrld,\ j \le n \cost{\pi},\ \tenv_2
        \text{ with } \dom(\tenv_2) \disjoint \chset
        \text{ and } \mathit{ok}_{\wrld'}(\senv_1 , \ch : \sAt{\sty}{\poff{\pntA}{1}}, \tenv_2),\
        \val \in \lrVW{\wrld'}{j}{\tenv_2 \types \payloadAt{\payty}{\pntA}} \st
      \\
      (\chset' \cup B, \pool') \in
        \lrPW{\wrld'}{j}{(\senv_1 , \ch : \sAt{\sty}{\poff{\pntA}{1}}) \mathbin{\triangleright} \tenv_2 \types \ty}
    }{
      \obl{\wrld ; \senv_1 , \ch : \sAt{(\sRecv{\payty}{\sty})}{\pntA}}{n}{(\chset,\pool) \trans{\lRecv{\ch}{\val}} (\chset',\pool') : \ty}
    } \label{rule:Q-Recv}
    \and
    \infer*[lab=Q-Send]{
      \dom(\tenv_{\asTy{e}}) = \fchans(\val) \cap \chset \\
      \dom(\tenv_{\asTy{e}}) \disjoint \dom(\senv_0 , \senv_{\asTy{v}}) \cup \set{\ch} \\
      \tenv_{\asTy{e}} \text{ and } \senv_{\asTy{e}} \text{ are dual at dual points}
      \\\\
      \val \in \lrVW{\wrld}{n}{\senv_{\asTy{v}} , \tenv_{\asTy{e}} \types \payloadAt{\payty}{\pntA}}
      \quad
      (\chset',\pool') \in \lrPW{\wrld}{n \cost{\pi}}{\senv_0 , \senv_{\asTy{e}} , \ch : \sAt{\sty}{\poff{\pntA}{1}} \types \ty}
    }{
      \obl{\wrld ; \senv_0 , \senv_{\asTy{v}} , \ch : \sAt{(\sSend{\payty}{\sty})}{\pntA}}{n}{(\chset,\pool) \trans{\lSend{\ch}{\val}} (\chset',\pool') : \ty}
    } \label{rule:Q-Send}
    \and
    \infer*[lab=Q-Close]{
      (\chset',\pool') \in \lrPW{\wrld}{n \cost{\pi}}{\senv_1 \types \ty}
    }{
      \obl{\wrld ; \senv_1 , \ch : \sAt{\sTerm}{\pntA}}{n}{(\chset,\pool) \trans{\lTerm{\ch}} (\chset',\pool') : \ty}
    } \label{rule:Q-Close}
    \and
    \infer*[lab=Q-Link]{
      \ch \ne \chB \\
      \sty_{\chB} \isdual \sty_{\ch} \\
      (\chset',\pool') \in \lrPW{\wrld}{n \cost{\pi}}{\senv_1 \types \ty}
    }{
      \obl{\wrld ; \senv_1 , \ch : \sAt{\sty_{\ch}}{\pntA} , \chB : \sAt{\sty_{\chB}}{\pdual{\pntA}}}{n}{(\chset,\pool) \trans{\lLink{\ch}{\chB}} (\chset',\pool') : \ty}
    } \label{rule:Q-Link}
  \end{proofrules}
  \rulesectionend
  \caption{The process logical relation}
  \Description{The two rules of the process relation and the five
    obligation rules of a process transition.}
  \label{fig:logrel-rules-proc}
\end{figure}

\Cref{rule:Q-Tau} is the only clause that grows the world, by the fresh
extension that \ref{rule:P-Act-New} needs.
For the $\tau$ of \ref{rule:P-Act-Link}, we require $N = \emptyset$
and $\wrld' = \wrld$ at index $n$.
At strict worlds, this single premise is equivalent to \rocq{Pobl}
by \rocq{wext_fresh_refl} and \rocq{Prel_down}.
In \ref{rule:Q-Recv} the message may bring the other end of a channel
the receiver already holds, or both ends of a fresh one. Those channels
$B$ become restricted in the continuation and their entries leave both
environments, which the operation $\senv \mathbin{\triangleright} \tenv_2$
performs (\rocq{recv_back}, \rocq{recv_env}). The premise
$\mathit{ok}_{\wrld'}(\senv, \tenv_2)$ asks that every single end in
$\tenv_2$ be dual to the end of its channel in $\senv$ at a dual
point, that the two ends of a channel arriving whole be dual at dual
points, and that such a channel be absent from $\senv$
(\rocq{recv_ok}).
\Cref{rule:Q-Send} is where a private channel becomes part of the
external interface. The restricted channels the message carries are
$\dom(\tenv_{\asTy{e}}) = \fchans(\val) \cap \chset$; the ends that
leave with the message are typed in $\tenv_{\asTy{e}}$, the ends that
stay in $\senv_{\asTy{e}}$, and the two are dual at dual points
(\rocq{send_split}). Two ends are at dual points when one sits at
$\pntA$ and the other at $\pdual{\pntA}$. The continuation then serves
$\senv_{\asTy{e}}$ as part of its interface. No condition on the
topology of the continuation is imposed: the order on points does that
work, so this rule needs no split of the continuation into one process
per extruded channel.

The syntactic side condition $\fchans(\val)$ in \ref{rule:Q-Send} is
again a consequence of the design of the LTS (\cref{rm:sentc}).

\subsection{Fundamental Property}
\label{sec:fundamental}

The term relation applies to closed terms, so an open term is closed by
a grounding substitution.

\begin{definition}[Grounding substitutions \rocq{grounding}]
\label{def:grounding}
  \[
    \lrGW{\wrld}{n}{\tenv ; \vr_1 : \ty_1, \dots, \vr_k : \ty_k} :=
      \left\{
        [\vr_1 \mapsto \val_1,\dots,\vr_k \mapsto \val_k] \;\middle|\;
          \begin{matrix}
          \tenv = \tenv_1 \cdot {} \dots \cdot \tenv_k
            \\
          \forall i.\,
            \val_i \in \lrVW{\wrld}{n}{\tenv_i \types \ty_i}
          \end{matrix}
      \right\}
  \]
\end{definition}

\begin{definition}[Semantic judgments \rocq{sem_typed}, \rocq{sem_ptyped}]
\label{def:sem-judgments}
  The judgment $\tenv ; \lenv \provesF{\frm}_{\wrld} \tm : \ty$ holds if
  \[
    \forall n \in \Nat.\,
      \forall \wrld' \mathrel{\sqsupseteq} \wrld.\,
        \forall \tenv'', \gamma \in \lrGW{\wrld'}{n}{\tenv' ; \lenv}
          \text{ with } \tenv'' = \tenv \cdot \tenv'.\,
          \tm\gamma \in \lrEW{\wrld'}{n}{\tenv'' \types \ty}{\frm}
  \]
  The judgment $\senv \proves_{\wrld} \proc : \ty$ holds if
  $
    \forall n \in \Nat.\,
      \proc \in \lrPW{\wrld}{n}{\senv \types \ty}
  $.
\end{definition}

The judgment is purely semantic. It says nothing about the syntax of
$\tm$ beyond what the term relation says about its behaviour. The two
syntactic conditions of \cref{rm:vabs} live in a separate predicate
$\mathit{wf}(\tenv; \lenv, \tm)$ (\rocq{term_wf}). It asks that the free
channels of $\tm$ be $\dom(\tenv)$, that every free variable of $\tm$
be bound in $\lenv$ and that $\tm$ mention every variable of $\lenv$ at
a type that is not unlimited. It holds of every
syntactically typed term (\rocq{typed_term_wf_s}) and enters as a
premise of exactly the two compositionality lemmas that build a
suspended computation.

\begin{lemma}[\rocq{Prel_steps}]
\label{lm:lrproc-inv}
  If $\proc \in \lrPW{\wrld}{n+m}{\tEmpty \types \ty}$ and
  $\proc \redu[n] \proc'$, then
  $\proc' \in \lrPW{\wrld'}{m}{\tEmpty \types \ty}$ for some $\wrld'$.
\end{lemma}

(\rocq{Prel_S_tau}).

\begin{lemma}[\rocq{Prel_safe}]
\label{lm:lrproc-safe}
  If $\proc \in \lrPW{\wrld}{n+1}{\tEmpty \types \ty}$ then $\proc$ is
  safe.
\end{lemma}

(\rocq{proc_progress}, \rocq{Prel_S_irred}, \rocq{Prel_S_empty_no_sess}).

\begin{theorem}[Adequacy \rocq{adequacy}]
  \label{thm:adequacy}
  If $\tEmpty \proves_{\wrld} \proc : \ty$ then $\proc$ is deadlock
  free.
\end{theorem}

\begin{proof}
  Let $\proc \redu[m] \proc'$. Pick $n = m+1$ in
  $\tEmpty \proves_{\wrld} \proc : \ty$. By
  \cref{lm:lrproc-inv}, $\proc' \in \lrPW{\wrld'}{1}{\tEmpty \types \ty}$
  for some $\wrld'$, and by \cref{lm:lrproc-safe} $\proc'$ is safe.
\end{proof}

The relation validates every rule of the static system. Read each rule
of \cref{fig:typeSys,fig:typeSysPoints,fig:procTypeSys} as an
implication, with $\provesF{\frm}_{\wrld}$ for the term rules and
$\proves_{\wrld}$ for the process rules, and the result is a valid
statement about the semantic judgment, with the two syntactic premises
noted above. The term rules are validated for the world-carrying
judgment of \rocq{typed}, which \rocq{styped_typed} reaches from a
source derivation.
\Cref{tab:compositionality} names the declaration for each rule, and the
three lemmas below spell out the rules that carry a side condition of
their own.

\begin{table}[tp]
  \caption{Compositionality: one declaration per rule of the static
    system.}
  \label{tab:compositionality}
  \begin{tabular}{llllll}
    \toprule
    Rule & Declaration & Rule & Declaration & Rule & Declaration \\
    \midrule
    \ref{rule:E-Unit} & \rocq{sem_typed_unit}
      & \ref{rule:E-Const} & \rocq{sem_typed_konst}
      & \ref{rule:E-Var} & \rocq{sem_typed_var} \\
    \ref{rule:E-Abs} & \rocq{sem_typed_lam}
      & \ref{rule:E-App} & \rocq{sem_typed_app}
      & \ref{rule:E-Pair} & \rocq{sem_typed_pair} \\
    \ref{rule:E-LetPair} & \rocq{sem_typed_letp}
      & \ref{rule:E-Let*} & \rocq{sem_typed_letu}
      & \ref{rule:E-Inj} & \rocq{sem_typed_inj1} \\
    \ref{rule:E-Case} & \rocq{sem_typed_case}
      & \ref{rule:E-Weak} & \rocq{sem_typed_wk}
      & \ref{rule:E-Contract} & \rocq{sem_typed_contract} \\
    \ref{rule:E-Conv} & \rocq{sem_typed_teq}
      & \ref{rule:E-Exp} & \rocq{sem_typed_box}
      & \ref{rule:E-Pack} & \rocq{sem_typed_pack} \\
    \ref{rule:E-Open} & \rocq{sem_typed_open}
      & \ref{rule:P-Term} & \rocq{Prel_term}
      & \ref{rule:P-Connect} & \rocq{Prel_connect} \\
    \bottomrule
  \end{tabular}
\end{table}

A channel constant is typed by its own rule in the world-carrying
judgment, with declaration \rocq{sem_typed_chan}, and
$\tApp{\tForce}{\tm}$ is \ref{rule:E-App} at $\tForce$
(\rocq{sem_typed_force}). Composing two processes that do not interact
is \ref{rule:P-Connect} with no shared channel (\rocq{Prel_indep}).

\begin{lemma}[Compositionality: \ref{rule:E-Abs} \rocq{sem_typed_lam}]
  \label{lem:comp:abs}
  If $\tenv ; \lenv , \vr : \ty \provesF{\frm'}_{\wrld'} \tm : \tyB$ for
  every $\wrld' \mathrel{\sqsupseteq} \wrld$ and every $\frm'$ with
  $\wrld' \types \fLt{\ubd}{\frm'}$,
  and $\wrld \types \fLe{\lbd}{\heads(\tenv) \cup \heads(\lenv)}$ and
  $\mathit{wf}(\tenv ; \lenv , \vr : \ty, \tm)$, then
  $\tenv ; \lenv \provesF{\frm}_{\wrld} \tAbs{\vr}{\tm}
    : \ty \lolli{\ubd}{\lbd} \tyB$.
\end{lemma}

(\rocq{Vrel_Erel}, \rocq{fchans_subst}, \rocq{grounding_heads}).

\begin{lemma}[Compositionality: \ref{rule:E-Exp} \rocq{sem_typed_box}]
  \label{lem:comp:exp}
  If
  $\tEmpty ; \lenv , \vr : \tyBox{\pvivec}{\cstr}{\tyUnit \lolli{\ubd}{} \ty}
     \provesF{\frm'}_{\wrld'} \tm : \tysub{\ty}{\pntvec}{\pvivec}$
  for every $\wrld' \mathrel{\sqsupseteq} \wrld$, every $\pntvec$ known to
  $\wrld'$ with $\wrld' \types \tysub{\cstr}{\pntvec}{\pvivec}$ and every
  $\frm'$ with $\wrld' \types \fLt{\tysub{\ubd}{\pntvec}{\pvivec}}{\frm'}$,
  with $\lenv$ of $\Box$ types only and
  $\mathit{wf}(\tEmpty ; \lenv , \vr : \tyBox{\pvivec}{\cstr}{\tyUnit \lolli{\ubd}{} \ty}, \tm)$,
  then
  $\tEmpty ; \lenv \provesF{\frm}_{\wrld} \tBoxP{\vr}{\tm}
    : \tyBox{\pvivec}{\cstr}{\tyUnit \lolli{\ubd}{} \ty}$.
\end{lemma}

(\rocq{grounding_unl}, \rocq{fchans_subst}, \rocq{Vrel_box_thunk_iff}).

\begin{lemma}[Compositionality: \ref{rule:P-Connect} \rocq{conn_all}, \rocq{Prel_connect}]
  \label{lem:comp:connect}
  If $\primitive \ty$, the two interfaces
  $\senv_{\asTy{C}}$ and $\dual{\senv_{\asTy{C}}}$ are dual at dual
  points of $\wrld$,
  $(\chset_1, \pool_1) \in \lrPW{\wrld}{n}{\senv_1 , \senv_{\asTy{C}} \types \ty}$,
  $(\chset_2, \pool_2) \in \lrPW{\wrld}{n}{\senv_2 , \dual{\senv_{\asTy{C}}} \types \tyUnit}$,
  the restricted channels of each component avoid the interface of the
  other, and $\wrld$ is strict, then
  \[
    (\chset_1 \cupdot \chset_2 \cupdot \dom(\senv_{\asTy{C}}),
     \pool_1 \uplus \pool_2)
      \in \lrPW{\wrld}{n}{\senv_1 , \senv_2 \types \ty}
  \]
\end{lemma}

\begin{proof}
  \rocq{conn_all} proves the indexed statement in one strict world.
  Monotonicity lifts both components to $\wrld_1\sqcup\wrld_2$
  (\rocq{Prel_mono}, \rocq{wext_join_l}).
  \Cref{thm:fresh-extension} preserves strictness when either component
  creates a channel (\rocq{wext_fresh_join}). This gives
  \rocq{Prel_connect}. Folding this cut gives the finite compositions
  (\rocq{Prel_connect_list}, \rocq{Prel_connect_n}).
\end{proof}

(\rocq{conn_all}, \rocq{Prel_connect}, \rocq{Prel_mono}, \rocq{wext_join_l},
\rocq{wext_fresh_join}, \rocq{cut_ok}, \rocq{ctys}, \rocq{ctys_sym},
\rocq{cut_ok_sym}, \rocq{conn_at}, \rocq{conn_lt}, \rocq{msize},
\rocq{pcost_giving}, \rocq{conn_level}, \rocq{Prel_S_intro}, \rocq{Prel_S_stop_prim},
\rocq{Vrel_prim_empty}, \rocq{conn_stop}, \rocq{link_l}, \rocq{link_r},
\rocq{stable_l}, \rocq{stable_r}, \rocq{ready_core}, \rocq{Prel_S_ready},
\rocq{cut_not_both}, \rocq{sdual_snd}, \rocq{sdual_rcv}, \rocq{sdual_end},
\rocq{msg_tau}, \rocq{close_tau}, \rocq{ext_l}, \rocq{step_l},
\rocq{step_r}, \rocq{descent}, \rocq{no_offers_l}, \rocq{no_offers_r},
\rocq{empty_of_no_offers}, \rocq{ready_connect}, \rocq{ptau_inv}, \rocq{proc_step_sess_inv},
\rocq{own_tau}, \rocq{Prel_tau_rename}, \rocq{tau_beta}, \rocq{tau_new},
\rocq{tau_spawn}, \rocq{tau_msg}, \rocq{tau_close}, \rocq{cross_msg},
\rocq{cross_close}, \rocq{cross_subj}, \rocq{ext_recv}, \rocq{ext_send},
\rocq{ext_close}, \rocq{ext_link}, \rocq{link_give}, \rocq{link_own},
\rocq{ext_link_norm}, \rocq{Prel_connect_list}, \rocq{Prel_connect_n}).

\Cref{lem:comp:connect} is the linear-logic cut, and it is the one place
where strictness of the composed world is needed:
\cref{thm:amalgamation,thm:fresh-extension} supply it along the run,
because the only clause that grows a world grows it freshly
(\ref{rule:Q-Tau}) and the names it adds are unknown to the other
component.

Partner points are what make the chain argument for readiness of the
composite go through (\rocq{ready_connect}, \rocq{descent}). In a stable
composite the two ends of a cut channel are never both offered. Take an
offered end at $\pntA$ on a cut channel whose partner end, at
$\pdual{\pntA}$, is unoffered. The component that holds the partner end
offers some entry at a point $\pntB$ with
$\wrld \types \fLt{\pntB}{\pdual{\pdual{\pntA}}}$, that is
$\fLt{\pntB}{\pntA}$. The offered points descend strictly, so the chain
visits each cut channel at most once and ends at an external offer. Read
as a deadlock argument, the blocked end with the least point $\pntA$ has
a partner at $\pdual{\pntA}$. If that partner is blocked at $\pntB$, it
owes $\fLt{\pntB}{\pntA}$, against the choice of $\pntA$. Otherwise it
offers the dual action and the two threads step.

\begin{theorem}[Fundamental property \rocq{fundamental}, \rocq{fundamental_proc}]
\label{th:fundamental}
  Typing implies semantic typing, at every strict world. The fact context
  of the term case is the world alone.
  \begin{itemize}
    \item
      $
        \tenv ; \wrld ; \lenv \typesF{\frm} \tm : \ty
        \implies
        \tenv ; \lenv \provesF{\frm}_{\wrld} \tm : \ty
      $
    \item
      $
        \senv \typesW{\wrld} \proc : \ty
        \implies
        \senv \proves_{\wrld} \proc : \ty
      $
  \end{itemize}
\end{theorem}

\begin{proof}
  Induction on the typing derivation, with one compositionality lemma
  per rule. \Cref{tab:compositionality} names the declaration for each.

  The term case runs in two stages. One induction on
  $\tenv ; \fctx ; \lenv \typesF{\frm} \tm : \ty$ instantiates the future
  variables and substitutes points, and lands in an infinitary judgment
  with one constructor per rule of the static system
  (\rocq{typed_sem_rules}, \rocq{sem_rules}). The constructors for
  $\lambda$ and $\Box$ carry their bodies at every later world, under
  every future above the upper point, and for a $\Box$ at every
  admissible instance of its constraint, which is what the two
  compositionality lemmas above ask for. A second induction on that
  judgment applies one lemma per constructor
  (\rocq{sem_rules_fundamental}). A closed derivation has no future
  variable and no hypothesis, so the identity substitution discharges
  the two stages' side conditions (\rocq{typed_sem_rules_closed}).

  The process case is induction on $\senv \typesW{\wrld} \proc : \ty$ with two
  rules. \Cref{rule:P-Term} is \rocq{Prel_term} applied to the term case,
  and \ref{rule:P-Connect} is \rocq{Prel_connect} at
  $\wrld_1 = \wrld_2 = \wrld$, which lands in $\wrld \sqcup \wrld$ and
  comes back to $\wrld$ by monotonicity, because $\wrld \sqcup \wrld$
  lies below $\wrld$ (\rocq{wext_join_self}, \rocq{Prel_mono}).
\end{proof}

\begin{theorem}[Soundness \rocq{soundness}]
\label{th:soundness}
  If $\wrld$ is strict and $\tEmpty \typesW{\wrld} \proc : \ty$, then
  $\proc$ is deadlock free.
\end{theorem}

\begin{proof}
  \Cref{th:fundamental} gives $\tEmpty \proves_{\wrld} \proc : \ty$, and
  \cref{thm:adequacy} turns that into deadlock freedom.
\end{proof}

\Cref{thm:soundness-source} is \cref{th:soundness} for a source program.
A source derivation reads no ground fact, the creation constraint
$\cstr_0$ supplies a strict world that realizes it, and the derivation
moves to that world (\rocq{styped_typed}).

\begin{theorem}[Soundness for source programs \rocq{soundness_source}]
\label{thm:soundness-source}
  If $\cstr_0$ is a creation constraint against the empty context,
  $\primitive \ty$ and $\cstr_0 ; \tEmpty \typesF{\emptyset} \tm : \ty$,
  then the thread $(\emptyset, \mset{\tm})$ is deadlock free.
\end{theorem}

\begin{proof}
  The proof builds the world the run needs out of $\cstr_0$. Embed the
  source derivation into the empty world, which supports no fact beyond
  program order (\rocq{wnil}, \rocq{wnil_strict},
  \rocq{styped_typed}). A creation constraint against the empty context
  is acyclic (\rocq{c0_creation_ok}). Let $n$ be the
  number of top-level point variables and send each of them to its own
  point constant. Acyclicity of $\cstr_0$ makes the world generated by
  the image of $\cstr_0$ under that map strict (\rocq{wc0},
  \rocq{wc0_strict}), it extends the empty world (\rocq{wc0_ext}), and
  the hypothesis $(\cstr_0, \pvivec)$ holds in it at those constants
  (\rocq{wc0_holds}). Instantiating the derivation along the map
  therefore yields the infinitary judgment in that world
  (\rocq{typed_sem_rules}, \rocq{c0_consts}), hence a semantic judgment
  (\rocq{sem_rules_fundamental}). One thread with nothing restricted is
  a process of the relation (\rocq{Prel_term}), and
  \cref{thm:adequacy} concludes. Point instantiation preserves
  $\primitive$ (\rocq{prim_pinst}), which \rocq{Prel_term} needs.
\end{proof}

%%% Local Variables:
%%% mode: latex
%%% TeX-master: "main"
%%% End:
